% \iffalse meta-comment
%
% Copyright (C) 2026 by Andrew Langan (alangan444@icloud.com)
%
% \fi
%
% \iffalse
%<*driver>
\ProvidesFile{os31lib.dtx}
\documentclass[hyperindex=false]{ltxdoc}

\usepackage{hypdoc}
\usepackage{bookmark}
\usepackage{dtxdescribe}
\usepackage{lmodern}
\usepackage{makeidx}
\usepackage{os31lib}

\EnableCrossrefs
\RecordChanges
\PageIndex

\begin{document}
  \DocInput{os31lib.dtx}
  \PrintChanges
  \PrintIndex
\end{document}
%</driver>
% \fi
%
% \changes{v0.0.4}{2026/10/27}{Bundled CWL and JSON files with package}
% \changes{v0.0.3}{2026/10/27}{Added Example Links}
% \changes{v0.0.2}{2026/10/27}{Added CWL and JSON files}
% \changes{v0.0.1}{2026/10/26}{Converted to DTX file}
%
% \DoNotIndex{\(,\),\[,\],\\,\{,\},\&,\%,\$,\#,\_,\^,\~,\,}
% \DoNotIndex{
%   \anchorX,\anchorY,\alph,\arabic,\BODY,
%   \begin,\DeclarePairedDelimiterX,\def,\delimsize,
%   \detokenize,\draw,\else,\end,\errmessage,
%   \ExplSyntaxOff,\ExplSyntaxOn,\fi,\fill,
%   \forall,\if,\ifdim,\ifinsol,\IfNoValueF,\IfNoValueT,
%   \in,\infty,\insolfalse,\insoltrue,\int,
%   \langle,\ldots,\left,\lim,\limits,\lvert,
%   \mathbb,\mathopen,\medskip,\newcounter,\NewEnviron,
%   \newif,\newpage,\newtheorem,\nexists,\node,
%   \normalfont,\notin,\omega,\par,\range,\relax,
%   \right,\rightarrow,\roman,\rvert,\section,
%   \setcounter,\small,\space,\stepcounter,\sum,
%   \text,\textbf,\theoremstyle,\vert,\ysign
% }
% \DoNotIndex{
%   \newcommand,
%   \newenvironment,
%   \RequirePackage,
%   \DeclarePairedDelimiter
% }
%
% \title{The \textsf{os31lib} package}
% \author{Andrew Langan}
% \date{\today}
% \maketitle
%
% \section{Introduction}
%
% This package is designed to make writing homework assignments in LaTeX easier.
% It provides a number of useful macros and environments for formatting
% problems, solutions, and mathematical expressions.
%
% \section{Usage}
%
% To use this package, include the following line in your preamble:
% \begin{verbatim}
%   \usepackage{os31lib}
% \end{verbatim}
%
% \newpage
% \section{Description}
%
% \subsection{Quantum Computing}
%
% \begin{description}
% \ItemDescribeMacro{\bra}{\marg{value}}
%
% Creates a bra vector in Dirac notation:
% \[
%   \langle \text{value} \|
% \]
%
% \ItemDescribeMacro{\ket}{\marg{value}}
%
% Creates a ket vector in Dirac notation:
% \[
%   \| \text{value} \rangle
% \]
%
% \ItemDescribeArgument[\bra, \ket]{value}
% The value to be enclosed in the bra or ket vector.
%
% \ItemDescribeMacro{\braket}{\marg{left} \marg{right}}
%
% Creates a bra-ket vector in Dirac notation:
% \[
%   \langle \text{left} \| \text{right} \rangle
% \]
%
% \ItemDescribeArgument[\braket]{left}
% The value to be enclosed in the bra vector.
%
% \ItemDescribeArgument[\braket]{right}
% The value to be enclosed in the ket vector.
% \end{description}
%
% \begin{description}
% \ItemDescribeMacro{\cplot}{\marg{x-max} \marg{real} \marg{imaginary} \marg{scale}}
%
% Plots a complex number on the complex plane using TikZ.
%
% \ItemDescribeArgument[\cplot]{x-max}
% The maximum value of the x-axis.
%
% \ItemDescribeArgument[\cplot]{real}
% The real part of the complex number.
%
% \ItemDescribeArgument[\cplot]{imaginary}
% The imaginary part of the complex number.
%
% \ItemDescribeArgument[\cplot]{scale}
% The scale of the plot.
%
% \end{description}
%
%
% See \hyperref[ex:Quantum Computing]{example\space \ref*{ex:Quantum Computing}}
%   for a demonstration of the above macros.
%
% \newpage
% \subsection{Enclosed Blocks}
%
% \begin{description}
% \ItemDescribeMacro{\bl}{\marg{block}}
%
% Creates a block of text enclosed in curly braces.
%
% \ItemDescribeMacro{\br}{\marg{block}}
%
% Creates a block of text enclosed in square brackets.
%
% \ItemDescribeMacro{\pr}{\marg{block}}
%
% Creates a block of text enclosed in parentheses.
%
% \ItemDescribeMacro{\abs}{\marg{block}}
%
% Creates a block of text enclosed in absolute value bars.
%
% \ItemDescribeArgument[\bl, \br, \pr, \abs]{block}
% The block of text to be enclosed.
% \end{description}
%
%
% See \hyperref[ex:Enclosed Blocks]{example\space \ref*{ex:Enclosed Blocks}} for
%   a demonstration of the above macros.
%
% \subsection{Comparison Operators}
%
% \begin{description}
% \ItemDescribeMacro{\eq}
%
% Represents the equality symbol: `\(=\)'.
%
% \ItemDescribeMacro{\gt}
%
% Represents the greater than symbol: `\(>\)'.
%
% \ItemDescribeMacro{\lt}
%
% Represents the less than symbol: `\(<\)'.
%
% \ItemDescribeMacro{\minus}
%
% Represents the minus symbol: `\(-\)'.
%
% \ItemDescribeMacro{\plus}
%
% Represents the plus symbol: `\(+\)'.
% \end{description}
%
%
% See \hyperref[ex:Comparison Operators]{example\space \ref*{ex:Comparison Operators}}
%   for a demonstration of the above macros.
%
% \newpage
% \subsection{Number Sets}
%
% \begin{description}
% \ItemDescribeMacro{\C}
%
% Represents the set of complex numbers: `\(\mathbb{C}\)'.
%
% \ItemDescribeMacro{\N}
%
% Represents the set of natural numbers: `\(\mathbb{N}\)'.
%
% \ItemDescribeMacro{\Q}
%
% Represents the set of rational numbers: `\(\mathbb{Q}\)'.
%
% \ItemDescribeMacro{\R}
%
% Represents the set of real numbers: `\(\mathbb{R}\)'.
%
% \ItemDescribeMacro{\W}
%
% Represents the set of whole numbers: `\(\omega\)'.
%
% \ItemDescribeMacro{\Z}
%
% Represents the set of integers: `\(\mathbb{Z}\)'.
% \end{description}
%
%
% See \hyperref[ex:Number Sets]{example\space \ref*{ex:Number Sets}} for a
%   demonstration of the above macros.
%
% \subsection{Let Statements}
%
% \begin{description}
% \ItemDescribeMacro{\LET}{\marg{variable} \marg{value}}
%
% Represents a let statement in math mode.
%
% \ItemDescribeMacro{\LETA}{\marg{variable} \marg{value}}
%
% Represents a let statement in math mode with alignment.
%
% \ItemDescribeArgument[\LET, \LETA]{variable}
% The variable being defined in the let statement.
%
% \ItemDescribeArgument[\LET, \LETA]{value}
% The value being assigned to the variable in the let statement.
% \end{description}
%
% See \hyperref[ex:Let Statements]{example\space \ref*{ex:Let Statements}} for a
%   demonstration of the above macros.
%
% \newpage
% \subsection{Set Membership}
%
% \begin{description}
% \ItemDescribeMacro{\inn}{\marg{variable}}
%
% Represents a variable in the set of natural numbers:
% \[
%   \text{variable} \in \mathbb{N}
% \]
%
% \ItemDescribeMacro{\ip}{\marg{variable}}
%
% Represents a variable in the set of prime numbers:
% \[
%   \text{variable} \in \mathbb{P}
% \]
%
% \ItemDescribeMacro{\iq}{\marg{variable}}
%
% Represents a variable in the set of rational numbers:
% \[
%   \text{variable} \in \mathbb{Q}
% \]
%
% \ItemDescribeMacro{\ir}{\marg{variable}}
%
% Represents a variable in the set of real numbers:
% \[
%   \text{variable} \in \mathbb{R}
% \]
%
% \ItemDescribeMacro{\iw}{\marg{variable}}
%
% Represents a variable in the set of whole numbers:
% \[
%   \text{variable} \in \omega
% \]
%
% \ItemDescribeMacro{\iz}{\marg{variable}}
%
% Represents a variable in the set of integers:
% \[
%   \text{variable} \in \mathbb{Z}
% \]
%
% \ItemDescribeArgument[\inn, \ip, \iq, \ir, \iw, \iz]{variable}
% The variable being checked for membership in the specified set.
% \end{description}
%
%
% See \hyperref[ex:Set (Non-)Membership]{example\space \ref*{ex:Set (Non-)Membership}}
%   for a demonstration of the above macros.
%
% \newpage
% \subsection{Set Non-Membership}
%
% \begin{description}
% \ItemDescribeMacro{\nin}{\marg{variable}}
%
% Represents a variable not in the set of natural numbers:
% \[
%   \text{variable} \notin \mathbb{N}
% \]
%
% \ItemDescribeMacro{\nip}{\marg{variable}}
%
% Represents a variable not in the set of prime numbers:
% \[
%   \text{variable} \notin \mathbb{P}
% \]
%
% \ItemDescribeMacro{\niq}{\marg{variable}}
%
% Represents a variable not in the set of rational numbers:
% \[
%   \text{variable} \notin \mathbb{Q}
% \]
%
% \ItemDescribeMacro{\nir}{\marg{variable}}
%
% Represents a variable not in the set of real numbers:
% \[
%   \text{variable} \notin \mathbb{R}
% \]
%
% \ItemDescribeMacro{\niw}{\marg{variable}}
%
% Represents a variable not in the set of whole numbers:
% \[
%   \text{variable} \notin \omega
% \]
%
% \ItemDescribeMacro{\niz}{\marg{variable}}
%
% Represents a variable not in the set of integers:
% \[
%   \text{variable} \notin \mathbb{Z}
% \]
%
% \ItemDescribeArgument[\nin, \nip, \niq, \nir, \niw, \niz]{variable}
% The variable being checked for non-membership in the specified set.
% \end{description}
%
%
% See \hyperref[ex:Set (Non-)Membership]{example\space \ref*{ex:Set (Non-)Membership}}
%   for a demonstration of the above macros.
%
% \newpage
% \subsection{Existential Quantifiers}
%
% \begin{description}
% \ItemDescribeMacro{\ein}{\marg{variable}}
%
% Represents the existential quantifier for a variable in the set of natural
% numbers:
% \[
%   (\exists \text{ variable} \in \mathbb{N})
% \]
%
% \ItemDescribeMacro{\eip}{\marg{variable}}
%
% Represents the existential quantifier for a variable in the set of prime
% numbers:
% \[
%   (\exists \text{ variable} \in \mathbb{P})
% \]
%
% \ItemDescribeMacro{\eiq}{\marg{variable}}
%
% Represents the existential quantifier for a variable in the set of rational
% numbers:
% \[
%   (\exists \text{ variable} \in \mathbb{Q})
% \]
%
% \ItemDescribeMacro{\eir}{\marg{variable}}
%
% Represents the existential quantifier for a variable in the set of real
% numbers:
% \[
%   (\exists \text{ variable} \in \mathbb{R})
% \]
%
% \ItemDescribeMacro{\eiw}{\marg{variable}}
%
% Represents the existential quantifier for a variable in the set of whole
% numbers:
% \[
%   (\exists \text{ variable} \in \omega)
% \]
%
% \ItemDescribeMacro{\eiz}{\marg{variable}}
%
% Represents the existential quantifier for a variable in the set of integers:
% \[
%   (\exists \text{ variable} \in \mathbb{Z})
% \]
%
% \ItemDescribeArgument[\ein, \eip, \eiq, \eir, \eiw, \eiz]{variable}
% The variable being quantified in the specified set.
% \end{description}
%
%
% See \hyperref[ex:Existential Quantifiers]{example\space \ref*{ex:Existential Quantifiers}}
%   for a demonstration of the above macros.
%
% \newpage
% \subsection{Existential Quantifiers for Non-Membership}
% \begin{description}
% \ItemDescribeMacro{\nein}{\marg{variable}}
%
% Represents the existential quantifier for a variable not in the set of natural
% numbers:
% \[
%   (\exists \text{ variable} \notin \mathbb{N})
% \]
%
% \ItemDescribeMacro{\neip}{\marg{variable}}
%
% Represents the existential quantifier for a variable not in the set of prime
% numbers:
% \[
%   (\exists \text{ variable} \notin \mathbb{P})
% \]
%
% \ItemDescribeMacro{\neiq}{\marg{variable}}
%
% Represents the existential quantifier for a variable not in the set of
% rational numbers:
% \[
%   (\exists \text{ variable} \notin \mathbb{Q})
% \]
%
% \ItemDescribeMacro{\neir}{\marg{variable}}
%
% Represents the existential quantifier for a variable not in the set of real
% numbers:
% \[
%   (\exists \text{ variable} \notin \mathbb{R})
% \]
%
% \ItemDescribeMacro{\neiw}{\marg{variable}}
%
% Represents the existential quantifier for a variable not in the set of whole
% numbers:
% \[
%   (\exists \text{ variable} \notin \omega)
% \]
%
% \ItemDescribeMacro{\neiz}{\marg{variable}}
%
% Represents the existential quantifier for a variable not in the set of
% integers:
% \[
%   (\exists \text{ variable} \notin \mathbb{Z})
% \]
%
% \ItemDescribeArgument[\nein, \neip, \neiq, \neir, \neiw, \neiz]{variable}
% The variable being quantified in the specified set.
% \end{description}
%
%
% See \hyperref[ex:Existential Quantifiers]{example\space \ref*{ex:Existential Quantifiers}}
%   for a demonstration of the above macros.
%
% \newpage
% \subsection{Universal Quantifiers}
%
% \begin{description}
% \ItemDescribeMacro{\ain}{\marg{variable}}
%
% Represents the universal quantifier for a variable in the set of natural
% numbers:
% \[
%   (\forall \text{ variable} \in \mathbb{N})
% \]
%
% \ItemDescribeMacro{\aip}{\marg{variable}}
%
% Represents the universal quantifier for a variable in the set of prime
% numbers:
% \[
%   (\forall \text{ variable} \in \mathbb{P})
% \]
%
% \ItemDescribeMacro{\aiq}{\marg{variable}}
%
% Represents the universal quantifier for a variable in the set of rational
% numbers:
% \[
%   (\forall \text{ variable} \in \mathbb{Q})
% \]
%
% \ItemDescribeMacro{\air}{\marg{variable}}
%
% Represents the universal quantifier for a variable in the set of real numbers:
% \[
%   (\forall \text{ variable} \in \mathbb{R})
% \]
% \ItemDescribeMacro{\aiw}{\marg{variable}}
%
% Represents the universal quantifier for a variable in the set of whole
% numbers:
% \[
%   (\forall \text{ variable} \in \omega)
% \]
% \ItemDescribeMacro{\aiz}{\marg{variable}}
%
% Represents the universal quantifier for a variable in the set of integers:
% \[
%   (\forall \text{ variable} \in \mathbb{Z})
% \]
%
% \ItemDescribeMacro{\nain}{\marg{variable}}
%
% Represents the universal quantifier for a variable not in the set of natural
% numbers:
% \[
%   (\forall \text{ variable} \notin \mathbb{N})
% \]
%
% \ItemDescribeArgument[\nain, \naiq, \nair, \naiw, \naiz]{variable}
% The variable being quantified in the specified set.
% \end{description}
%
%
% See \hyperref[ex:Universal Quantifiers]{example\space \ref*{ex:Universal Quantifiers}}
%   for a demonstration of the above macros.
%
% \newpage
% \subsection{Infinite Sums and Integrals}
%
% \begin{description}
% \ItemDescribeMacro{\sumk}
%
% Represents the summation from \(k = 1\) to \(n\):
% \[
%   \sum\limits_{k = 1}^{n} \ldots
% \]
%
% \ItemDescribeMacro{\sumkz}
%
% Represents the summation from \(k = 0\) to \(n\):
% \[
%   \sum\limits_{k = 0}^{n} \ldots
% \]
%
% \ItemDescribeMacro{\sumn}
%
% Represents the summation from \(i = 1\) to \(n\):
% \[
%   \sum\limits_{i = 1}^{n} \ldots
% \]
%
% \end{description}
%
% \begin{description}
% \ItemDescribeMacro{\infint}{\marg{integrand}}
%
% Represents the integral from 0 to \(\infty\):
% \[
%   \int_{0}^{\infty} \text{integrand} \, dt
% \]
%
% \ItemDescribeArgument[\infint]{integrand}
% The integrand of the integral.
% \end{description}
%
% \begin{description}
% \ItemDescribeMacro{\inflim}{\marg{block}}
%
% Represents the limit as \(x\) approaches infinity:
% \[
%   \lim_{n \to \infty}( \text{block} )
% \]
%
% \ItemDescribeArgument[\inflim]{block}
% The block of text to be evaluated in the limit.
% \end{description}
%
%
% See \hyperref[ex:Infinite Sums and Integrals]{example\space \ref*{ex:Infinite Sums and Integrals}}
%   for a demonstration of the above macros.
%
% \newpage
% \subsection{Telescoping Series}
%
% \begin{description}
% \ItemDescribeMacro{\ts}{\marg{n}}
%
% Represents the summation of a telescoping series:
% \[
%   1^{n} + 2^{n} + \ldots + n^{n}
% \]
%
% \ItemDescribeMacro{\ets}{\marg{n}}
%
% Represents the summation of an extended telescoping series:
% \[
%   1^{n} - 0^{n} + 2^{n} - 1^{n} + 3^{n} - 2^{n} + \ldots + n^{n} - (n-1)^{n}
% \]
%
% \ItemDescribeArgument[\ts, \ets]{n}
% The exponent of the telescoping series.
% \end{description}
%
%
% See \hyperref[ex:Telescoping Series]{example\space \ref*{ex:Telescoping Series}}
%   for a demonstration of the above macros.
%
% \subsection{Miscellaneous}
%
% \begin{description}
% \ItemDescribeMacro{\nl}
%
% Inserts a new line in the document.
% \end{description}
%
%
% See \hyperref[ex:Miscellaneous]{example\space \ref*{ex:Miscellaneous}} for a
%   demonstration of the above macros.
%
% \newpage
% \subsection{Environments}
%
% \begin{description}
% \ItemDescribeEnv{prob}{\oarg{title}}
%
% Environment for formatting a homework question. It automatically numbers the
% problems, given an optional title.
%
% \ItemDescribeArgument[\prob]{title}
% The optional title of the homework question.
%
% \ItemDescribeEnv{subprob}
%
% Environment for formatting a subproblem of a homework question. It
% automatically numbers the sub-problems.
%
% \ItemDescribeEnv{subsubprob}
%
% Environment for formatting a sub-subproblem of a homework question. It
% automatically numbers the sub-sub-problems.
%
% \ItemDescribeEnv{sol}
%
% Environment for formatting a solution to a homework question.
%
% \ItemDescribeEnv{indproof}{\marg{base case} \marg{assumption} \marg{inductive case}}
%
% Environment for formatting an inductive proof.
%
% \ItemDescribeArgument[\indproof]{base case}
% The base case of the inductive proof.
%
% \ItemDescribeArgument[\indproof]{assumption}
% The assumption of the inductive proof.
%
% \ItemDescribeArgument[\indproof]{inductive case}
% The inductive case of the inductive proof.
%
% \ItemDescribeEnv{page}{\oarg{title}}
%
% Environment for formatting a page of a homework assignment. Takes up the
% entire page.
%
% \ItemDescribeArgument[\page]{title}
% The optional title of the homework assignment page.
% \end{description}
%
%
% See \hyperref[ex:Environments]{example\space \ref*{ex:Environments}} for a
%   demonstration of the above environments.
%
% \newpage
% \section{Implementation}
%
% \StopEventually{}
%
% \iffalse
%<*os31lib>
% \fi
%
%    \begin{macrocode}
\ExplSyntaxOn

\RequirePackage{amssymb, amsthm, amsmath, color}
\RequirePackage{array}
\RequirePackage{environ}
\RequirePackage{etoolbox}
\RequirePackage{float}
\RequirePackage{graphics}
\RequirePackage{graphicx}
\RequirePackage{mathtools}
\RequirePackage{subfigure}
\RequirePackage{tikz}

\theoremstyle{definition}
\newtheorem{problem}{Problem}
\newtheorem{theorem}{Theorem}[section]
\newtheorem{lemma}[theorem]{Lemma}
\newtheorem{proposition}[theorem]{Proposition}
\newtheorem{claim}[theorem]{Claim}
\newtheorem{corollary}[theorem]{Corollary}
\newtheorem{definition}[theorem]{Definition}
\newtheorem{observation}[theorem]{Observation}
\newtheorem{fact}[theorem]{Fact}
\newtheorem{property}{Property}
\newtheorem{remark}{Remark}[section]
\newtheorem{notation}{Notation}[section]
\newtheorem{example}{Example}[section]
\newtheorem{algorithm}{Algorithm}
\newtheorem{conjecture}{Conjecture}
\newtheorem{question}[conjecture]{Question}

\newif\ifinsol
\insolfalse

\newcounter{prob}
\newcounter{subprob}
\newcounter{subsubprob}
%    \end{macrocode}
%
% \iffalse
% region Quantum Computing
% \fi
%
%    \begin{macrocode}
\DeclarePairedDelimiter\bra{\langle}{\rvert}
\DeclarePairedDelimiter\ket{\lvert}{\rangle}
\DeclarePairedDelimiterX\braket[2]{\langle}{\rangle}{
  #1 \, \delimsize\vert \, \mathopen{} #2
} % chktex 21
%    \end{macrocode}
%
% \begin{macro}{\cplot} \marg{x-max} \marg{real} \marg{imaginary} \marg{scale}
%    \begin{macrocode}
\newcommand{\cplot}[4]{
  \begin{tikzpicture}[scale=#4]
    \IfNoValueT{#1}{
      \errmessage{The\space x-max\space of\space the\space complex\space
        plane\space is\space not\space provided.}
    }
    \IfNoValueT{#2}{
      \errmessage{The\space real\space part\space of\space the\space
        complex \space number \space is\space not\space provided.}
    }
    \IfNoValueT{#3}{
      \errmessage{The\space imaginary\space part\space
        of\space the\space complex\space number\space
        is\space not\space provided.}
    }
    \IfNoValueT{#4}{
      \errmessage{The\space scale\space of\space the\space complex\space
        plane\space is\space not\space provided.}
    }
    \draw[gray!30,very\space thin] (-#1, -#1) grid (#1, #1);
    \draw[->] (-#1, 0) -- (#1, 0) node[right] {\(Re(z)\)};
    \draw[->] (0, -#1) -- (0, #1) node[above] {\(Im(z)\)};
    \draw[->, thick, blue] (0, 0) -- (#2, #3);

    \fill (#2, #3) circle (2pt);

    \ifdim #2pt > 0pt
      \def\anchorX{right}
    \else
      \def\anchorX{left}
    \fi
    \ifdim #3pt > 0pt
      \def\anchorY{above}
    \else
      \def\anchorY{below}
    \fi
    \ifdim #3pt < 0pt
      \def\ysign{}
    \else
      \def\ysign{+}
    \fi

    \node[\anchorY\space \anchorX] at (#2, #3)
    {\(z = #2 \ysign #3 i\)};
  \end{tikzpicture}
}
%    \end{macrocode}
% \end{macro}

% \iffalse
% endregion
% \fi
%
% \iffalse
% region Enclosed Blocks
% \fi
%
% \begin{macro}{\bl} \marg{block}
%    \begin{macrocode}
\newcommand{\bl}[1]{
  \left\{ #1 \right\} % chktex 21
}
%    \end{macrocode}
% \end{macro}
%
%
% \begin{macro}{\br} \marg{block}
%    \begin{macrocode}
\newcommand{\br}[1]{
  \left[ #1 \right]
}
%    \end{macrocode}
% \end{macro}
%
%
% \begin{macro}{\pr} \marg{block}
%    \begin{macrocode}
\newcommand{\pr}[1]{
  \left( #1 \right) % chktex 37
}
%    \end{macrocode}
% \end{macro}
%
%
% \begin{macro}{\abs} \marg{block}
%    \begin{macrocode}
\newcommand{\abs}[1]{
  \left\lvert #1 \right\rvert
}
%    \end{macrocode}
% \end{macro}

% \iffalse
% endregion
% \fi
%
% \iffalse
% region Comparison Operators
% \fi
%
% \begin{macro}{\eq}
%    \begin{macrocode}
\newcommand{\eq}{=}
%    \end{macrocode}
% \end{macro}
%
%
% \begin{macro}{\gt}
%    \begin{macrocode}
\newcommand{\gt}{>}
%    \end{macrocode}
% \end{macro}
%
%
% \begin{macro}{\lt}
%    \begin{macrocode}
\newcommand{\lt}{<}
%    \end{macrocode}
% \end{macro}
%
%
% \begin{macro}{\minus}
%    \begin{macrocode}
\newcommand{\minus}{-}
%    \end{macrocode}
% \end{macro}
%
%
% \begin{macro}{\plus}
%    \begin{macrocode}
\newcommand{\plus}{+}
%    \end{macrocode}
% \end{macro}
%
% \iffalse
% endregion
% \fi
%
% \iffalse
% region Number Sets
% \fi
%
% \begin{macro}{\C}
%    \begin{macrocode}
\newcommand{\C}{\mathbb{C}}
%    \end{macrocode}
% \end{macro}
%
%
% \begin{macro}{\N}
%    \begin{macrocode}
\newcommand{\N}{\mathbb{N}}
%    \end{macrocode}
% \end{macro}
%
%
% \begin{macro}{\Q}
%    \begin{macrocode}
\newcommand{\Q}{\mathbb{Q}}
%    \end{macrocode}
% \end{macro}
%
%
% \begin{macro}{\R}
%    \begin{macrocode}
\newcommand{\R}{\mathbb{R}}
%    \end{macrocode}
% \end{macro}
%
%
% \begin{macro}{\W}
%    \begin{macrocode}
\newcommand{\W}{\omega}
%    \end{macrocode}
% \end{macro}
%
%
% \begin{macro}{\Z}
%    \begin{macrocode}
\newcommand{\Z}{\mathbb{Z}}
%    \end{macrocode}
% \end{macro}

% \iffalse
% endregion
% \fi
%
% \iffalse
% region Let Statements
% \fi
%
% \begin{macro}{\LET} \marg{variable} \marg{value}
%    \begin{macrocode}
\newcommand{\LET}[2]{
  \text{Let{\space}} #1 = #2
}
%    \end{macrocode}
% \end{macro}
%
%
% \begin{macro}{\LETA} \marg{variable} \marg{value}
%    \begin{macrocode}
\newcommand{\LETA}[2]{
  \text{Let{\space}} #1 & = #2
}
%    \end{macrocode}
% \end{macro}

% \iffalse
% endregion
% \fi
%
% \iffalse
% region Set Membership
% \fi
%
% \begin{macro}{\inn} \marg{variable}
%    \begin{macrocode}
\newcommand{\inn}[1]{
  #1 \in \N
}
%    \end{macrocode}
% \end{macro}
%
%
% \begin{macro}{\ip} \marg{variable}
%    \begin{macrocode}
\newcommand{\ip}[1]{
  #1 \in \mathbb{P}
}
%    \end{macrocode}
% \end{macro}
%
%
% \begin{macro}{\iq} \marg{variable}
%    \begin{macrocode}
\newcommand{\iq}[1]{
  #1 \in \Q
}
%    \end{macrocode}
% \end{macro}
%
%
% \begin{macro}{\ir} \marg{variable}
%    \begin{macrocode}
\newcommand{\ir}[1]{
  #1 \in \R
}
%    \end{macrocode}
% \end{macro}
%
%
% \begin{macro}{\iw} \marg{variable}
%    \begin{macrocode}
\newcommand{\iw}[1]{
  #1 \in \W
}
%    \end{macrocode}
% \end{macro}
%
%
% \begin{macro}{\iz} \marg{variable}
%    \begin{macrocode}
\newcommand{\iz}[1]{
  #1 \in \Z
}
%    \end{macrocode}
% \end{macro}
%
%
% \begin{macro}{\nin} \marg{variable}
%    \begin{macrocode}
\newcommand{\nin}[1]{
  #1 \notin \N
}
%    \end{macrocode}
% \end{macro}
%
%
% \begin{macro}{\nip} \marg{variable}
%    \begin{macrocode}
\newcommand{\nip}[1]{
  #1 \notin \mathbb{P}
}
%    \end{macrocode}
% \end{macro}
%
%
% \begin{macro}{\niq} \marg{variable}
%    \begin{macrocode}
\newcommand{\niq}[1]{
  #1 \notin \Q
}
%    \end{macrocode}
% \end{macro}
%
%
% \begin{macro}{\nir} \marg{variable}
%    \begin{macrocode}
\newcommand{\nir}[1]{
  #1 \notin \R
}
%    \end{macrocode}
% \end{macro}
%
%
% \begin{macro}{\niw} \marg{variable}
%    \begin{macrocode}
\newcommand{\niw}[1]{
  #1 \notin \W
}
%    \end{macrocode}
% \end{macro}
%
%
% \begin{macro}{\niz} \marg{variable}
%    \begin{macrocode}
\newcommand{\niz}[1]{
  #1 \notin \Z
}
%    \end{macrocode}
% \end{macro}

% \iffalse
% endregion
% \fi
%
% \iffalse
% region Quantifiers
% \fi
%
% \begin{macro}{\ein} \marg{variable}
%    \begin{macrocode}
\newcommand{\ein}[1]{
  (\inn{\exists #1})
}
%    \end{macrocode}
% \end{macro}
%
%
% \begin{macro}{\ain} \marg{variable}
%    \begin{macrocode}
\newcommand{\ain}[1]{
  (\inn{\forall #1})
}
%    \end{macrocode}
% \end{macro}
%
%
% \begin{macro}{\eip} \marg{variable}
%    \begin{macrocode}
\newcommand{\eip}[1]{
  (\ip{\exists #1})
}
%    \end{macrocode}
% \end{macro}
%
%
% \begin{macro}{\aip} \marg{variable}
%    \begin{macrocode}
\newcommand{\aip}[1]{
  (\ip{\forall #1})
}
%    \end{macrocode}
% \end{macro}
%
%
% \begin{macro}{\eiq} \marg{variable}
%    \begin{macrocode}
\newcommand{\eiq}[1]{
  (\iq{\exists #1})
}
%    \end{macrocode}
% \end{macro}
%
%
% \begin{macro}{\aiq} \marg{variable}
%    \begin{macrocode}
\newcommand{\aiq}[1]{
  (\iq{\forall #1})
}
%    \end{macrocode}
% \end{macro}
%
%
% \begin{macro}{\eir} \marg{variable}
%    \begin{macrocode}
\newcommand{\eir}[1]{
  (\ir{\exists #1})
}
%    \end{macrocode}
% \end{macro}
%
%
% \begin{macro}{\air} \marg{variable}
%    \begin{macrocode}
\newcommand{\air}[1]{
  (\ir{\forall #1})
}
%    \end{macrocode}
% \end{macro}
%
%
% \begin{macro}{\eiw} \marg{variable}
%    \begin{macrocode}
\newcommand{\eiw}[1]{
  (\iw{\exists #1})
}
%    \end{macrocode}
% \end{macro}
%
%
% \begin{macro}{\aiw} \marg{variable}
%    \begin{macrocode}
\newcommand{\aiw}[1]{
  (\iw{\forall #1})
}
%    \end{macrocode}
% \end{macro}
%
%
% \begin{macro}{\eiz} \marg{variable}
%    \begin{macrocode}
\newcommand{\eiz}[1]{
  (\iz{\exists #1})
}
%    \end{macrocode}
% \end{macro}
%
%
% \begin{macro}{\aiz} \marg{variable}
%    \begin{macrocode}
\newcommand{\aiz}[1]{
  (\iz{\forall #1})
}
%    \end{macrocode}
% \end{macro}
%
%
% \begin{macro}{\nain} \marg{variable}
%    \begin{macrocode}
\newcommand{\nain}[1]{
  (\inn{\forall #1})
}
%    \end{macrocode}
% \end{macro}
%
%
% \begin{macro}{\nein} \marg{variable}
%    \begin{macrocode}
\newcommand{\nein}[1]{
  (\inn{\nexists #1})
}
%    \end{macrocode}
% \end{macro}
%
%
% \begin{macro}{\neip} \marg{variable}
%    \begin{macrocode}
\newcommand{\neip}[1]{
  (\ip{\nexists #1})
}
%    \end{macrocode}
% \end{macro}
%
%
% \begin{macro}{\neiq} \marg{variable}
%    \begin{macrocode}
\newcommand{\neiq}[1]{
  (\iq{\nexists #1})
}
%    \end{macrocode}
% \end{macro}
%
%
% \begin{macro}{\neir} \marg{variable}
%    \begin{macrocode}
\newcommand{\neir}[1]{
  (\ir{\nexists #1})
}
%    \end{macrocode}
% \end{macro}
%
%
% \begin{macro}{\neiw} \marg{variable}
%    \begin{macrocode}
\newcommand{\neiw}[1]{
  (\iw{\nexists #1})
}
%    \end{macrocode}
% \end{macro}
%
%
% \begin{macro}{\neiz} \marg{variable}
%    \begin{macrocode}
\newcommand{\neiz}[1]{
  (\iz{\nexists #1})
}
%    \end{macrocode}
% \end{macro}
%
% \iffalse
% endregion
% \fi
%
% \iffalse
% region Infinite Sums and Integrals
% \fi
%
% \begin{macro}{\sumk}
%    \begin{macrocode}
\ExplSyntaxOff
\newcommand{\sumk}{
  \sum\limits_{k \eq 1}^{n}
}
\ExplSyntaxOn
%    \end{macrocode}
% \end{macro}
%
%
% \begin{macro}{\sumkz}
%    \begin{macrocode}
\ExplSyntaxOff
\newcommand{\sumkz}{
  \sum\limits_{k \eq 0}^{n}
}
\ExplSyntaxOn
%    \end{macrocode}
% \end{macro}
%
%
% \begin{macro}{\sumn}
%    \begin{macrocode}
\ExplSyntaxOff
\newcommand{\sumn}{
  \sum\limits_{i \eq 1}^{n}
}
\ExplSyntaxOn
%    \end{macrocode}
% \end{macro}
%
%
% \begin{macro}{\infint} \marg{integrand}
%    \begin{macrocode}
\ExplSyntaxOff
\newcommand{\infint}[1]{
  \int\limits_{0}^{\infty} #1 \, dt
}
\ExplSyntaxOn
%    \end{macrocode}
% \end{macro}
%
%
% \begin{macro}{\inflim} \marg{block}
%    \begin{macrocode}
\ExplSyntaxOff
\newcommand{\inflim}[1]{
  \lim_{n \rightarrow \infty} \left( #1 \right) % chktex 37
}
\ExplSyntaxOn
%    \end{macrocode}
% \end{macro}
%
% \iffalse
% endregion
% \fi
%
% \iffalse
% region Telescoping Sums
% \fi
%
% \begin{macro}{\ts} \marg{power}
%    \begin{macrocode}
\newcommand{\ts}[1]{
  \IfNoValueT{#1}{
    \errmessage{
      The\space exponent\space of\space the\space telescoping\space
        series\space is\space not\space provided.
    }
  }
  1^{#1} \plus 2^{#1} \plus \ldots \plus n^{#1}
}
%    \end{macrocode}
% \end{macro}
%
%
% \begin{macro}{\ets} \marg{power}
%    \begin{macrocode}
\newcommand{\ets}[1]{
  \IfNoValueT{#1}{
    \errmessage{
      The\space exponent\space of\space the\space extended\space
        telescoping\space series\space is\space not\space provided.
    }
  }
  1^{#1} \minus 0^{#1}
  \plus 2^{#1} \minus 1^{#1}
  \plus 3^{#1} \minus 2^{#1}
  \plus \ldots
  \plus n^{#1} \minus {(n \minus 1)}^{#1}
}
%    \end{macrocode}
% \end{macro}

% \iffalse
% endregion
% \fi
%
% \iffalse
% region Miscellaneous
% \fi
%
% \begin{macro}{\nl}
%    \begin{macrocode}
\newcommand{\nl}{\par\medskip} % chktex 21
%    \end{macrocode}
% \end{macro}

% \iffalse
% endregion
% \fi
%
% \iffalse
% region Environments
% \fi
%
% \begin{environment}{prob} \oarg{title}
%    \begin{macrocode}
\newenvironment{prob}[1][]{
  \newpage
  \stepcounter{prob}
  \setcounter{subprob}{0}
  \setcounter{subsubprob}{0}
  \section*{Problem \arabic{prob}
    \if\relax\detokenize{#1}\relax
    \else{\normalfont\small(#1)}
    \fi
  }
}{}
%    \end{macrocode}
% \end{environment}
%
%
% \begin{environment}{subprob}
%    \begin{macrocode}
\newenvironment{subprob}{
  \stepcounter{subprob}
  \setcounter{subsubprob}{0}
  \subsection*{\alph{subprob}}
}{}
%    \end{macrocode}
% \end{environment}
%
%
% \begin{environment}{subsubprob}
%    \begin{macrocode}
\newenvironment{subsubprob}{
  \stepcounter{subsubprob}
  \subsubsection*{\roman{subsubprob}}
}{}
%    \end{macrocode}
% \end{environment}
%
%
% \begin{environment}{sol}
%    \begin{macrocode}
\newenvironment{sol}{
  \insoltrue
  \begin{proof}
  }{
  \end{proof}
  \insolfalse
}
%    \end{macrocode}
% \end{environment}
%
%
% \begin{environment}{indproof} \marg{base case} \marg{assumption}
%   \marg{inductive case}
%    \begin{macrocode}
\NewEnviron{indproof}[3][]{
  \IfNoValueT{#1}{
    \errmessage{
      Base\space case\space not\space provided.
    }
  }
  \IfNoValueT{#2}{
    \errmessage{
      Inductive\space assumption\space not\space provided.
    }
  }
  \IfNoValueT{#3}{
    \errmessage{
      Inductive\space case\space not\space provided.
    }
  }
  \ifinsol
    \errmessage{
      The\space indproof\space environment\space cannot\space be\space
        used\space inside\space the\space sol\space environment.
    }
  \fi

  \begin{sol}
    \nl
    \textbf{Base{\space}Case:}
    \par\medskip
    #1
    \nl
    \textbf{Inductive{\space}Step:}
    \nl
    Assume{\space}#2
    \nl
    \BODY
    \nl
    \textbf{Conclusion:}
    \nl
    #3
    \nl
  \end{sol}
}
%    \end{macrocode}
% \end{environment}
%
%
% \begin{environment}{page} \oarg{title}
%    \begin{macrocode}
\newenvironment{page}[1][]{
  \newpage
  \IfNoValueF{#1}{
    \textbf{#1}
  }
  \nl
}{}
%    \end{macrocode}
% \end{environment}
%
% \iffalse
% endregion
% \fi
%
\ExplSyntaxOff
%
% \iffalse
%</os31lib>
% \fi
%
% \iffalse
%<*os31libcwl>
%% #include:amsmath
%% #include:amssymb
%% #include:amsthm
%% #include:array
%% #include:color
%% #include:environ
%% #include:etoolbox
%% #include:float
%% #include:graphics
%% #include:graphicx
%% #include:mathtools
%% #include:subfigure
%% #include:tikz

%% # Quantum Computing
%% \bra{value%formula}#m
%% \ket{value%formula}#m
%% \braket{left%formula}{right%formula}#m
%% \cplot{x maximum}{real}{imaginary}{scale}#n

%% # Enclosed Blocks
%% \bl{block%formula}#m
%% \br{block%formula}#m
%% \pr{block%formula}#m
%% \abs{block%formula}#m

%% # Comparison Operators
%% \eq#m
%% \gt#m
%% \lt#m
%% \minus#m
%% \plus#m

%% # Number Sets
%% \C#m
%% \N#m
%% \Q#m
%% \R#m
%% \W#m
%% \Z#m

%% # Let Statements
%% \LET{variable name}{value}#m
%% \LETA{variable name}{value}#T

%% # Set Membership
%% \inn{variable name}#m
%% \ip{variable name}#m
%% \iq{variable name}#m
%% \ir{variable name}#m
%% \iw{variable name}#m
%% \iz{variable name}#m

%% # Set Non-Membership
%% \nin{variable name}#m
%% \nip{variable name}#m
%% \niq{variable name}#m
%% \nir{variable name}#m
%% \niw{variable name}#m
%% \niz{variable name}#m

%% # Existential Quantifiers
%% \ein{variable name}#m
%% \eip{variable name}#m
%% \eiq{variable name}#m
%% \eir{variable name}#m
%% \eiw{variable name}#m
%% \eiz{variable name}#m

%% # Existential Quantifiers for Non-Membership
%% \nein{variable name}#m
%% \neip{variable name}#m
%% \neiq{variable name}#m
%% \neir{variable name}#m
%% \neiw{variable name}#m
%% \neiz{variable name}#m

%% # Universal Quantifiers
%% \ain{variable name}#m
%% \aip{variable name}#m
%% \aiq{variable name}#m
%% \air{variable name}#m
%% \aiw{variable name}#m
%% \aiz{variable name}#m
%% \nain{variable name}#m

%% # Infinite Sums and Integrals
%% \sumk#m
%% \sumkz#m
%% \sumn#m
%% \infint{integrand%formula}#m
%% \inflim{block%formula}#m

%% # Telescoping Series
%% \ts{power%formula}#m
%% \ets{power%formula}#m

%% # Miscellaneous
%% \nl

%% # Environments
%% \begin{prob}[problem number]#beginEnv#prob
%% \begin{subprob}#beginEnv#subprob#/prob
%% \begin{subsubprob}#beginEnv#subsubprob#/subprob
%% \begin{sol}#beginEnv#sol#/prob,subprob,subsubprob
%% \begin{indproof}[base case]{inductive thesis%formula}{conclusion%formula}#beginEnv#indproof#/prob,subprob,subsubprob
%% \begin{page}[page title%title]

%% \end{prob}#endEnv#prob
%% \end{subprob}#endEnv#subprob
%% \end{subsubprob}#endEnv#subsubprob
%% \end{sol}#endEnv#sol
%% \end{indproof}#endEnv#indproof
%% \end{page}#endEnv#page
%</os31libcwl>
% \fi
%
%
% \iffalse
%<*os31libjson>
%% {
%%   "args": [],
%%   "deps": [
%%     {
%%       "name": "amsmath"
%%     },
%%     {
%%       "name": "amssymb"
%%     },
%%     {
%%       "name": "amsthm"
%%     },
%%     {
%%       "name": "array"
%%     },
%%     {
%%       "name": "color"
%%     },
%%     {
%%       "name": "environ"
%%     },
%%     {
%%       "name": "etoolbox"
%%     },
%%     {
%%       "name": "float"
%%     },
%%     {
%%       "name": "graphics"
%%     },
%%     {
%%       "name": "graphicx"
%%     },
%%     {
%%       "name": "mathtools"
%%     },
%%     {
%%       "name": "subfigure"
%%     },
%%     {
%%       "name": "tikz"
%%     }
%%   ],
%%   "envs": [
%%     {
%%       "arg": {
%%         "format": "[]",
%%         "snippet": "[${1:problem number}]"
%%       },
%%       "name": "prob"
%%     },
%%     {
%%       "name": "subprob"
%%     },
%%     {
%%       "name": "subsubprob"
%%     },
%%     {
%%       "name": "sol"
%%     },
%%     {
%%       "arg": {
%%         "format": "[]{}{}",
%%         "snippet": "[${3:base case}]{${1:inductive thesis%formula}}{${2:conclusion%formula}}"
%%       },
%%       "name": "indproof"
%%     },
%%     {
%%       "arg": {
%%         "format": "[]",
%%         "snippet": "[${1:page title%title}]"
%%       },
%%       "name": "page"
%%     }
%%   ],
%%   "keys": {},
%%   "macros": [
%%     {
%%       "arg": {
%%         "format": "{}",
%%         "snippet": "bra{${1:value%formula}}"
%%       },
%%       "detail": "Bra vector",
%%       "name": "bra"
%%     },
%%     {
%%       "arg": {
%%         "format": "{}",
%%         "snippet": "ket{${1:value%formula}}"
%%       },
%%       "detail": "Ket vector",
%%       "name": "ket"
%%     },
%%     {
%%       "arg": {
%%         "format": "{}{}",
%%         "snippet": "braket{${1:left%formula}}{${2:right%formula}}"
%%       },
%%       "detail": "Braket notation",
%%       "name": "braket"
%%     },
%%     {
%%       "arg": {
%%         "format": "{}{}{}{}",
%%         "snippet": "cplot{${1:x maximum}}{${2:real}}{${3:imaginary}}{${4:scale}}"
%%       },
%%       "detail": "Complex plot",
%%       "name": "cplot"
%%     },
%%     {
%%       "arg": {
%%         "format": "{}",
%%         "snippet": "bl{${1:block%formula}}"
%%       },
%%       "detail": "Braced block",
%%       "name": "bl"
%%     },
%%     {
%%       "arg": {
%%         "format": "{}",
%%         "snippet": "br{${1:block%formula}}"
%%       },
%%       "detail": "Bracketed block",
%%       "name": "br"
%%     },
%%     {
%%       "arg": {
%%         "format": "{}",
%%         "snippet": "pr{${1:block%formula}}"
%%       },
%%       "detail": "Parenthesized block",
%%       "name": "pr"
%%     },
%%     {
%%       "arg": {
%%         "format": "{}",
%%         "snippet": "abs{${1:block%formula}}"
%%       },
%%       "detail": "Absolute value",
%%       "name": "abs"
%%     },
%%     {
%%       "detail": "=",
%%       "name": "eq"
%%     },
%%     {
%%       "detail": ">",
%%       "name": "gt"
%%     },
%%     {
%%       "detail": "<",
%%       "name": "lt"
%%     },
%%     {
%%       "detail": "-",
%%       "name": "minus"
%%     },
%%     {
%%       "detail": "+",
%%       "name": "plus"
%%     },
%%     {
%%       "detail": "Complex numbers",
%%       "name": "C"
%%     },
%%     {
%%       "detail": "Natural numbers",
%%       "name": "N"
%%     },
%%     {
%%       "detail": "Rational numbers",
%%       "name": "Q"
%%     },
%%     {
%%       "detail": "Real numbers",
%%       "name": "R"
%%     },
%%     {
%%       "detail": "Whole numbers",
%%       "name": "W"
%%     },
%%     {
%%       "detail": "Integers",
%%       "name": "Z"
%%     },
%%     {
%%       "arg": {
%%         "format": "{}{}",
%%         "snippet": "LET{${1:variable name}}{${2:value}}"
%%       },
%%       "detail": "Let statement",
%%       "name": "LET"
%%     },
%%     {
%%       "arg": {
%%         "format": "{}{}",
%%         "snippet": "LETA{${1:variable name}}{${2:value}}"
%%       },
%%       "detail": "Let statement (aligned)",
%%       "name": "LETA"
%%     },
%%     {
%%       "arg": {
%%         "format": "{}",
%%         "snippet": "inn{${1:variable name}}"
%%       },
%%       "detail": "Within natural numbers",
%%       "name": "inn"
%%     },
%%     {
%%       "arg": {
%%         "format": "{}",
%%         "snippet": "ip{${1:variable name}}"
%%       },
%%       "detail": "Within prime numbers",
%%       "name": "ip"
%%     },
%%     {
%%       "arg": {
%%         "format": "{}",
%%         "snippet": "iq{${1:variable name}}"
%%       },
%%       "detail": "Within rational numbers",
%%       "name": "iq"
%%     },
%%     {
%%       "arg": {
%%         "format": "{}",
%%         "snippet": "ir{${1:variable name}}"
%%       },
%%       "detail": "Within real numbers",
%%       "name": "ir"
%%     },
%%     {
%%       "arg": {
%%         "format": "{}",
%%         "snippet": "iw{${1:variable name}}"
%%       },
%%       "detail": "Within whole numbers",
%%       "name": "iw"
%%     },
%%     {
%%       "arg": {
%%         "format": "{}",
%%         "snippet": "iz{${1:variable name}}"
%%       },
%%       "detail": "Within integers",
%%       "name": "iz"
%%     },
%%     {
%%       "arg": {
%%         "format": "{}",
%%         "snippet": "nin{${1:variable name}}"
%%       },
%%       "detail": "Not within natural numbers",
%%       "name": "nin"
%%     },
%%     {
%%       "arg": {
%%         "format": "{}",
%%         "snippet": "nip{${1:variable name}}"
%%       },
%%       "detail": "Not within prime numbers",
%%       "name": "nip"
%%     },
%%     {
%%       "arg": {
%%         "format": "{}",
%%         "snippet": "niq{${1:variable name}}"
%%       },
%%       "detail": "Not within rational numbers",
%%       "name": "niq"
%%     },
%%     {
%%       "arg": {
%%         "format": "{}",
%%         "snippet": "nir{${1:variable name}}"
%%       },
%%       "detail": "Not within real numbers",
%%       "name": "nir"
%%     },
%%     {
%%       "arg": {
%%         "format": "{}",
%%         "snippet": "niw{${1:variable name}}"
%%       },
%%       "detail": "Not within whole numbers",
%%       "name": "niw"
%%     },
%%     {
%%       "arg": {
%%         "format": "{}",
%%         "snippet": "niz{${1:variable name}}"
%%       },
%%       "detail": "Not within integers",
%%       "name": "niz"
%%     },
%%     {
%%       "arg": {
%%         "format": "{}",
%%         "snippet": "ein{${1:variable name}}"
%%       },
%%       "detail": "Exists within natural numbers",
%%       "name": "ein"
%%     },
%%     {
%%       "arg": {
%%         "format": "{}",
%%         "snippet": "eip{${1:variable name}}"
%%       },
%%       "detail": "Exists within prime numbers",
%%       "name": "eip"
%%     },
%%     {
%%       "arg": {
%%         "format": "{}",
%%         "snippet": "eiq{${1:variable name}}"
%%       },
%%       "detail": "Exists within rational numbers",
%%       "name": "eiq"
%%     },
%%     {
%%       "arg": {
%%         "format": "{}",
%%         "snippet": "eir{${1:variable name}}"
%%       },
%%       "detail": "Exists within real numbers",
%%       "name": "eir"
%%     },
%%     {
%%       "arg": {
%%         "format": "{}",
%%         "snippet": "eiw{${1:variable name}}"
%%       },
%%       "detail": "Exists within whole numbers",
%%       "name": "eiw"
%%     },
%%     {
%%       "arg": {
%%         "format": "{}",
%%         "snippet": "eiz{${1:variable name}}"
%%       },
%%       "detail": "Exists within integers",
%%       "name": "eiz"
%%     },
%%     {
%%       "arg": {
%%         "format": "{}",
%%         "snippet": "nein{${1:variable name}}"
%%       },
%%       "detail": "Does not exist within natural numbers",
%%       "name": "nein"
%%     },
%%     {
%%       "arg": {
%%         "format": "{}",
%%         "snippet": "neip{${1:variable name}}"
%%       },
%%       "detail": "Does not exist within prime numbers",
%%       "name": "neip"
%%     },
%%     {
%%       "arg": {
%%         "format": "{}",
%%         "snippet": "neiq{${1:variable name}}"
%%       },
%%       "detail": "Does not exist within rational numbers",
%%       "name": "neiq"
%%     },
%%     {
%%       "arg": {
%%         "format": "{}",
%%         "snippet": "neir{${1:variable name}}"
%%       },
%%       "detail": "Does not exist within real numbers",
%%       "name": "neir"
%%     },
%%     {
%%       "arg": {
%%         "format": "{}",
%%         "snippet": "neiw{${1:variable name}}"
%%       },
%%       "detail": "Does not exist within whole numbers",
%%       "name": "neiw"
%%     },
%%     {
%%       "arg": {
%%         "format": "{}",
%%         "snippet": "neiz{${1:variable name}}"
%%       },
%%       "detail": "Does not exist within integers",
%%       "name": "neiz"
%%     },
%%     {
%%       "arg": {
%%         "format": "{}",
%%         "snippet": "ain{${1:variable name}}"
%%       },
%%       "detail": "For all within natural numbers",
%%       "name": "ain"
%%     },
%%     {
%%       "arg": {
%%         "format": "{}",
%%         "snippet": "aip{${1:variable name}}"
%%       },
%%       "detail": "For all within prime numbers",
%%       "name": "aip"
%%     },
%%     {
%%       "arg": {
%%         "format": "{}",
%%         "snippet": "aiq{${1:variable name}}"
%%       },
%%       "detail": "For all within rational numbers",
%%       "name": "aiq"
%%     },
%%     {
%%       "arg": {
%%         "format": "{}",
%%         "snippet": "air{${1:variable name}}"
%%       },
%%       "detail": "For all within real numbers",
%%       "name": "air"
%%     },
%%     {
%%       "arg": {
%%         "format": "{}",
%%         "snippet": "aiw{${1:variable name}}"
%%       },
%%       "detail": "For all within whole numbers",
%%       "name": "aiw"
%%     },
%%     {
%%       "arg": {
%%         "format": "{}",
%%         "snippet": "aiz{${1:variable name}}"
%%       },
%%       "detail": "For all within integers",
%%       "name": "aiz"
%%     },
%%     {
%%       "arg": {
%%         "format": "{}",
%%         "snippet": "nain{${1:variable name}}"
%%       },
%%       "detail": "For all not within natural numbers",
%%       "name": "nain"
%%     },
%%     {
%%       "detail": "Summation from k = 1 to n",
%%       "name": "sumk"
%%     },
%%     {
%%       "detail": "Summation from k = 0 to n",
%%       "name": "sumkz"
%%     },
%%     {
%%       "detail": "Summation from i = 1 to n",
%%       "name": "sumn"
%%     },
%%     {
%%       "arg": {
%%         "format": "{}",
%%         "snippet": "infint{${1:integrand%formula}}"
%%       },
%%       "detail": "Integral from 0 to infinity",
%%       "name": "infint"
%%     },
%%     {
%%       "arg": {
%%         "format": "{}",
%%         "snippet": "inflim{${1:block%formula}}"
%%       },
%%       "detail": "Limit as x approaches infinity",
%%       "name": "inflim"
%%     },
%%     {
%%       "arg": {
%%         "format": "{}",
%%         "snippet": "ts{${1:power%formula}}"
%%       },
%%       "detail": "Telescoping series",
%%       "name": "ts"
%%     },
%%     {
%%       "arg": {
%%         "format": "{}",
%%         "snippet": "ets{${1:power%formula}}"
%%       },
%%       "detail": "Expanded telescoping series",
%%       "name": "ets"
%%     },
%%     {
%%       "detail": "New line",
%%       "name": "nl"
%%     }
%%   ]
%% }
%</os31libjson>
% \fi
%
%
% \iffalse
%<*letafix>
% \fi
\begin{align*}
  \LETA{y}{x + 2} \\
  y - x & = 2
\end{align*}
%
% \iffalse
%</letafix>
% \fi
%
% \Finale
%
% \iffalse
% Fix for the \texttt{LETA} macro, which does not render correctly in the
% documentation. This is a temporary fix until the issue is resolved.
% \fi
%
% \NewDocumentEnvironment{letafix}{s +O{} m}{}{
%   \par
%   \addvspace{\bigskipamount}
%   \IfBooleanTF{#1}{
%     \minipage{\linewidth}
%     \captionsetup{type=dtxexamplefloat}
%   }{
%     \begin{dtxexamplefloat}
%   }
%   \hrule\medskip
%   \caption{#3}
%   \textcolor{DTXD@examplerulecolor}{\smallskip\hrule}
%   \smallskip
%   {\scriptsize\itshape\dtxexamplecodename}
%   \VerbatimInput[tabsize=4]{leta_fix.tex}
%   \unskip
%   \textcolor{DTXD@examplerulecolor}{\hrule}
%   \smallskip
%   {\scriptsize\itshape\dtxexampleresultname}
%   \ifstrempty{#2}{}{
%     \itshape\small #2
%   }
%   \IfBooleanTF{#1}{
%     \endminipage
%   }{
%     \end{dtxexamplefloat}
%   }
% }
% \AfterEndEnvironment{letafix}{
%   \begin{align*}
%     \text{Let } y & = x + 2 \\
%     y - x & = 2
%   \end{align*}
%   \medskip\hrule
% }
%
% \newpage
% \section{Examples}
%
% \begin{dtxexample}
% {\env{Quantum Computing}\label{ex:Quantum Computing}}
% \begin{align*}
%   \bra{\psi} & = \bra{\psi_1} \otimes \bra{\psi_2} \\
%   \ket{\psi} & = \ket{\psi_1} \otimes \ket{\psi_2} \\
%   \braket{\psi}{\phi} & = \braket{\psi_1}{\phi_1}
%     \cdot \braket{\psi_2}{\phi_2}
% \end{align*}
%
% \cplot{10}{2}{3}{0.25}
% \cplot{10}{-2}{-3}{0.25}
% \end{dtxexample}
%
% \newpage
% \begin{dtxexample}
% {\env{Enclosed Blocks}\label{ex:Enclosed Blocks}}
% \begin{align*}
%  \bl{\frac{1 + 2 + 3}{2}} & = \bl{3} \\
%  \br{\frac{1 + 2 + 3}{2}} & = \br{3} \\
%  \pr{\frac{1 + 2 + 3}{2}} & = \pr{3} \\
%  \abs{\frac{1 + 2 + 3}{2}} & = \abs{3}
% \end{align*}
% \end{dtxexample}
%
% \begin{dtxexample}
% {\env{Comparison Operators}\label{ex:Comparison Operators}}
% \begin{align*}
%  1 \eq 1 & \\
%  1 \gt 0 & \\
%  0 \lt 1 & \\
%  1 \minus 1 & = 0 \\
%  1 \plus 1 & = 2
% \end{align*}
% \end{dtxexample}
%
% \newpage
% \begin{dtxexample}
% {\env{Number Sets}\label{ex:Number Sets}}
% \begin{align*}
%   \C & = \{ a + bi \mid a, b \in \R \} \\
%   \N & = \{ 1, 2, 3, \ldots \} \\
%   \Q & = \{ \frac{a}{b} \mid a, b \in \Z, b \neq 0 \} \\
%   \R & = \{ x \mid x \text{ is a real number} \} \\
%   \W & = \{ 0, 1, 2, \ldots \} \\
%   \Z & = \{ \ldots, -2, -1, 0, 1, 2, \ldots \}
% \end{align*}
% \end{dtxexample}
%
% \iffalse
% This example does not render correctly, but it is valid.
% \fi
% \begin{letafix}
% {\env{Let Statements}\label{ex:Let Statements}}
% \begin{align*}
%   \text{Let } y & = x + 2 \\
%   y - x & = 2
% \end{align*}
% \end{letafix}
%
% \newpage
% \begin{dtxexample}
% {\env{Set (Non-)Membership}\label{ex:Set (Non-)Membership}}
% \[
%   \begin{array}{|c|c|}
%     \hline
%     \inn{n} & \nin{n} \\
%     \ip{n} & \nip{n} \\
%     \iq{n} & \niq{n} \\
%     \ir{n} & \nir{n} \\
%     \iz{n} & \niz{n} \\
%     \iw{n} & \niw{n} \\
%     \hline
%   \end{array}
% \]
% \end{dtxexample}
%
% \newpage
% \begin{dtxexample}
% {\env{Existential Quantifiers}\label{ex:Existential Quantifiers}}
% \[
%   \begin{array}{|c|c|}
%     \hline
%     \ein{n} & \nein{n} \\
%     \eip{n} & \neip{n} \\
%     \eiq{n} & \neiq{n} \\
%     \eir{n} & \neir{n} \\
%     \eiz{n} & \neiz{n} \\
%     \eiw{n} & \neiw{n} \\
%     \hline
%   \end{array}
% \]
% \end{dtxexample}
%
% \begin{dtxexample}
% {\env{Universal Quantifiers}\label{ex:Universal Quantifiers}}
% \[
%   \begin{array}{|c|c|}
%     \hline
%     \ain{n} & \nain{n} \\
%     \aip{n} & \\
%     \aiq{n} & \\
%     \air{n} & \\
%     \aiw{n} & \\
%     \aiz{n} & \\
%     \hline
%   \end{array}
% \]
% \end{dtxexample}
%
% \begin{dtxexample}
% {\env{Infinite Sums and Integrals}\label{ex:Infinite Sums and Integrals}}
% \begin{align*}
%   \sumk 1 & = n \\
%   \sumkz 1 & = n + 1 \\
%   \sumn 1 & = n \\
%   \infint{e^{-t}} & = 1 \\
%   \inflim{n} \frac{1}{n} & = 0
% \end{align*}
% \end{dtxexample}
%
% \begin{dtxexample}
% {\env{Telescoping Series}\label{ex:Telescoping Series}}
% \begin{align*}
%   \sumk k^2 & = \ts{2} \\
%   \sumk (k^2 - (k-1)^2) & = \ets{2}
% \end{align*}
% \end{dtxexample}
%
% \newpage
% \begin{dtxexample}
% {\env{Miscellaneous}\label{ex:Miscellaneous}}
% Line 1
% (Not Line 2)
% \nl
% Line 2
% \nl
% Line 3
% \end{dtxexample}
%
% \begin{dtxexample}
% {\env{Environments}\label{ex:Environments}}
% \newpage
% \begin{problem}[1.1]
%   \nl
%   Solve the following:
%   \begin{subprob}
%    Show by induction that \(n^2 + n\) is even for all \(n \in \N\).
%
%     \begin{subsubprob}
%       Show the base case for \(n = 1\).
%
%       \begin{sol}
%         \begin{align*}
%           1^2 + 1 & = 1 + 1 \\
%           & = 2 \\
%           & \text{is even.}
%         \end{align*}
%       \end{sol}
%     \end{subsubprob}
%
%     \begin{subsubprob}
%       Show the full proof.
%
%       \begin{indproof}[\(n = 1\)]
%         {\(n^2 + n\) is even for \(n \in \N\).}
%         {Since \(n^2 + n\) is even by the inductive hypothesis,
%           and \(2(n + 1)\) is even, the sum is even.}
%         \begin{align*}
%           (n + 1)^2 + (n + 1) & = n^2 + 2n + 1 + n + 1 \\
%           & = n^2 + n + 2n + 2 \\
%           & = (n^2 + n) + 2(n + 1)
%         \end{align*}
%       \end{indproof}
%     \end{subsubprob}
%   \end{subprob}
% \end{problem}
% \end{dtxexample}
%
\endinput
%
